%%%%%%%%%%%%%%%%%%%%%%%%%%%%%%%%%%%%%%%%%%%%%%%%%%%%%%%%%%%%
% 6.2 BASIC COUNTING
% The Composition of Advanced Mathematics
%%%%%%%%%%%%%%%%%%%%%%%%%%%%%%%%%%%%%%%%%%%%%%%%%%%%%%%%%%%%

\section{Basic Counting}
\label{sec:basic-counting}

\prerequisite{
Finite sets, Cartesian products, functions, subsets, factorial notation,
basic algebra, and the discrete structures introduced in
Section~\ref{sec:discrete-structures}.
}

\learningobjectives{
By the end of this section, the reader should be able to:
\begin{itemize}
    \item apply the addition and multiplication principles;
    \item use complements to simplify counting problems;
    \item apply the division principle;
    \item use the pigeonhole principle and its generalized form;
    \item distinguish permutations, arrangements, and combinations;
    \item count selections with and without repetition;
    \item evaluate and manipulate binomial coefficients;
    \item apply Pascal's identity and the binomial theorem;
    \item count permutations with repeated objects;
    \item count ordered and unordered distributions;
    \item use stars and bars;
    \item apply inclusion--exclusion to finite sets;
    \item count surjective functions;
    \item recognize derangements;
    \item use double counting to prove combinatorial identities;
    \item construct bijective counting arguments;
    \item distinguish exact counting from overcounting and undercounting;
    \item translate verbal counting problems into mathematical models.
\end{itemize}
}

%%%%%%%%%%%%%%%%%%%%%%%%%%%%%%%%%%%%%%%%%%%%%%%%%%%%%%%%%%%%
\subsection{The Fundamental Question of Counting}
\label{subsec:fundamental-question-counting}
%%%%%%%%%%%%%%%%%%%%%%%%%%%%%%%%%%%%%%%%%%%%%%%%%%%%%%%%%%%%

Combinatorics begins with a simple question:

\[
\boxed{
\text{How many objects satisfy the required conditions?}
}
\]

The challenge is rarely arithmetic alone.

The main difficulty is identifying:

\begin{itemize}
    \item what constitutes one outcome;
    \item which outcomes are distinct;
    \item whether order matters;
    \item whether repetition is allowed;
    \item whether restrictions are present;
    \item whether multiple descriptions represent the same object.
\end{itemize}

A correct counting argument therefore begins by clearly identifying the
set being counted.

If \(S\) is the set of admissible outcomes, the problem is to determine

\[
|S|.
\]

%%%%%%%%%%%%%%%%%%%%%%%%%%%%%%%%%%%%%%%%%%%%%%%%%%%%%%%%%%%%
\subsection{The Addition Principle}
\label{subsec:addition-principle}
%%%%%%%%%%%%%%%%%%%%%%%%%%%%%%%%%%%%%%%%%%%%%%%%%%%%%%%%%%%%

Suppose a collection of outcomes is divided into disjoint cases.

\begin{theorem}[Addition Principle]
If finite sets \(A_1,\ldots,A_k\) are pairwise disjoint, then

\[
\boxed{
\left|
\bigcup_{i=1}^{k}A_i
\right|
=
\sum_{i=1}^{k}|A_i|.
}
\]
\end{theorem}

The phrase \emph{pairwise disjoint} means

\[
A_i\cap A_j=\varnothing
\qquad
(i\neq j).
\]

Thus if an outcome can occur in exactly one of several cases, the
numbers of possibilities are added.

%%%%%%%%%%%%%%%%%%%%%%%%%%%%%%%%%%%%%%%%%%%%%%%%%%%%%%%%%%%%
\subsection{Worked Example: Addition Principle}
\label{subsec:worked-addition-principle}
%%%%%%%%%%%%%%%%%%%%%%%%%%%%%%%%%%%%%%%%%%%%%%%%%%%%%%%%%%%%

\begin{workedexamplebox}{Selecting One Course}

Suppose a student may choose one course from either:

\[
8
\]

mathematics courses or

\[
5
\]

computer science courses.

Assume the two lists contain distinct courses.

The choices fall into two disjoint classes.

Therefore,

\[
8+5=13.
\]

\begin{answerbox}
\[
\boxed{13}
\]

courses are available.
\end{answerbox}

\end{workedexamplebox}

%%%%%%%%%%%%%%%%%%%%%%%%%%%%%%%%%%%%%%%%%%%%%%%%%%%%%%%%%%%%
\subsection{The Multiplication Principle}
\label{subsec:multiplication-principle}
%%%%%%%%%%%%%%%%%%%%%%%%%%%%%%%%%%%%%%%%%%%%%%%%%%%%%%%%%%%%

Many counting problems consist of a sequence of choices.

\begin{theorem}[Multiplication Principle]
Suppose a process consists of \(k\) successive choices. If the first
choice has \(n_1\) possibilities, and after each previous choice the
\(i\)-th choice always has \(n_i\) possibilities, then the total number
of outcomes is

\[
\boxed{
n_1n_2\cdots n_k.
}
\]
\end{theorem}

This principle corresponds to the cardinality of a Cartesian product.

If

\[
S=A_1\times A_2\times\cdots\times A_k,
\]

then

\[
\boxed{
|S|
=
\prod_{i=1}^{k}|A_i|.
}
\]

%%%%%%%%%%%%%%%%%%%%%%%%%%%%%%%%%%%%%%%%%%%%%%%%%%%%%%%%%%%%
\subsection{Worked Example: Multiplication Principle}
\label{subsec:worked-multiplication-principle}
%%%%%%%%%%%%%%%%%%%%%%%%%%%%%%%%%%%%%%%%%%%%%%%%%%%%%%%%%%%%

\begin{workedexamplebox}{Creating an Identification Code}

Suppose a code contains:

\begin{itemize}
    \item two letters;
    \item followed by three decimal digits.
\end{itemize}

Assume repetition is allowed.

Each letter position has

\[
26
\]

choices.

Each digit position has

\[
10
\]

choices.

Thus

\[
26\cdot26\cdot10\cdot10\cdot10
=
26^2 10^3.
\]

\begin{answerbox}
\[
\boxed{
676000
}
\]

different codes are possible.
\end{answerbox}

\end{workedexamplebox}

%%%%%%%%%%%%%%%%%%%%%%%%%%%%%%%%%%%%%%%%%%%%%%%%%%%%%%%%%%%%
\subsection{Dependent Choices}
\label{subsec:dependent-choices}
%%%%%%%%%%%%%%%%%%%%%%%%%%%%%%%%%%%%%%%%%%%%%%%%%%%%%%%%%%%%

The multiplication principle does not require the same number of
options at every step.

Suppose the first choice has

\[
n
\]

possibilities.

After making that choice, the second has

\[
n-1
\]

possibilities.

The third then has

\[
n-2
\]

possibilities.

The total is

\[
\boxed{
n(n-1)(n-2).
}
\]

This pattern appears whenever objects cannot be reused.

%%%%%%%%%%%%%%%%%%%%%%%%%%%%%%%%%%%%%%%%%%%%%%%%%%%%%%%%%%%%
\subsection{Decision Trees}
\label{subsec:counting-decision-trees}
%%%%%%%%%%%%%%%%%%%%%%%%%%%%%%%%%%%%%%%%%%%%%%%%%%%%%%%%%%%%

A counting process can often be visualized as a tree.

Each level represents a choice.

Each branch represents one possible decision.

Every root-to-leaf path represents one complete outcome.

If every node at level \(i\) has the same number \(n_i\) of outgoing
branches, the number of leaves is

\[
\boxed{
n_1n_2\cdots n_k.
}
\]

Decision trees provide a visual interpretation of the multiplication
principle.

%%%%%%%%%%%%%%%%%%%%%%%%%%%%%%%%%%%%%%%%%%%%%%%%%%%%%%%%%%%%
\subsection{The Complement Principle}
\label{subsec:complement-principle}
%%%%%%%%%%%%%%%%%%%%%%%%%%%%%%%%%%%%%%%%%%%%%%%%%%%%%%%%%%%%

Sometimes it is easier to count what is forbidden than what is allowed.

If \(A\subseteq U\), then

\[
|A^c|
=
|U|-|A|.
\]

Equivalently,

\[
\boxed{
|A|
=
|U|-|A^c|.
}
\]

This is called counting by complement.

%%%%%%%%%%%%%%%%%%%%%%%%%%%%%%%%%%%%%%%%%%%%%%%%%%%%%%%%%%%%
\subsection{Worked Example: Counting by Complement}
\label{subsec:worked-counting-complement}
%%%%%%%%%%%%%%%%%%%%%%%%%%%%%%%%%%%%%%%%%%%%%%%%%%%%%%%%%%%%

\begin{workedexamplebox}{Binary Strings Containing at Least One \(1\)}

How many binary strings of length \(8\) contain at least one \(1\)?

There are

\[
2^8
\]

binary strings of length \(8\).

The only string containing no \(1\)'s is

\[
00000000.
\]

Therefore,

\[
2^8-1=255.
\]

\begin{answerbox}
\[
\boxed{255}
\]

binary strings contain at least one \(1\).
\end{answerbox}

\end{workedexamplebox}

%%%%%%%%%%%%%%%%%%%%%%%%%%%%%%%%%%%%%%%%%%%%%%%%%%%%%%%%%%%%
\subsection{The Division Principle}
\label{subsec:division-principle}
%%%%%%%%%%%%%%%%%%%%%%%%%%%%%%%%%%%%%%%%%%%%%%%%%%%%%%%%%%%%

Sometimes a construction counts each desired object the same number of
times.

\begin{theorem}[Division Principle]
If a counting procedure produces \(N\) representations and every
underlying object is represented exactly \(k\) times, then the number
of distinct objects is

\[
\boxed{
\frac{N}{k}.
}
\]
\end{theorem}

The critical condition is that every object must be counted exactly the
same number of times.

%%%%%%%%%%%%%%%%%%%%%%%%%%%%%%%%%%%%%%%%%%%%%%%%%%%%%%%%%%%%
\subsection{Worked Example: Handshakes}
\label{subsec:worked-handshakes}
%%%%%%%%%%%%%%%%%%%%%%%%%%%%%%%%%%%%%%%%%%%%%%%%%%%%%%%%%%%%

\begin{workedexamplebox}{Counting Pairs of People}

Suppose \(n\) people are present.

One way to count ordered choices of two distinct people is

\[
n(n-1).
\]

However, the pair

\[
(A,B)
\]

and

\[
(B,A)
\]

represent the same handshake.

Thus every unordered pair is counted twice.

By the division principle,

\[
\frac{n(n-1)}{2}.
\]

\begin{answerbox}
\[
\boxed{
\frac{n(n-1)}{2}
}
\]

handshakes are possible if every pair shakes hands once.
\end{answerbox}

\end{workedexamplebox}

%%%%%%%%%%%%%%%%%%%%%%%%%%%%%%%%%%%%%%%%%%%%%%%%%%%%%%%%%%%%
\subsection{The Pigeonhole Principle}
\label{subsec:pigeonhole-principle}
%%%%%%%%%%%%%%%%%%%%%%%%%%%%%%%%%%%%%%%%%%%%%%%%%%%%%%%%%%%%

The pigeonhole principle converts a counting inequality into an
existence statement.

\begin{theorem}[Pigeonhole Principle]
If more than \(n\) objects are placed into \(n\) boxes, then at least
one box contains at least two objects.
\end{theorem}

In function language, if

\[
f:A\to B
\]

with

\[
|A|>|B|,
\]

then \(f\) cannot be injective.

%%%%%%%%%%%%%%%%%%%%%%%%%%%%%%%%%%%%%%%%%%%%%%%%%%%%%%%%%%%%
\subsection{Worked Example: Shared Birth Month}
\label{subsec:worked-pigeonhole}
%%%%%%%%%%%%%%%%%%%%%%%%%%%%%%%%%%%%%%%%%%%%%%%%%%%%%%%%%%%%

\begin{workedexamplebox}{Birth Months}

Suppose \(13\) people are present.

There are only

\[
12
\]

possible birth months.

Treat the people as objects and the months as boxes.

Since

\[
13>12,
\]

at least two people must have birthdays in the same month.

\begin{answerbox}
\[
\boxed{
\text{At least two people share a birth month.}
}
\]
\end{answerbox}

\end{workedexamplebox}

%%%%%%%%%%%%%%%%%%%%%%%%%%%%%%%%%%%%%%%%%%%%%%%%%%%%%%%%%%%%
\subsection{Generalized Pigeonhole Principle}
\label{subsec:generalized-pigeonhole}
%%%%%%%%%%%%%%%%%%%%%%%%%%%%%%%%%%%%%%%%%%%%%%%%%%%%%%%%%%%%

The pigeonhole principle can be strengthened.

\begin{theorem}[Generalized Pigeonhole Principle]
If \(N\) objects are distributed among \(k\) boxes, then some box
contains at least

\[
\boxed{
\left\lceil \frac{N}{k}\right\rceil
}
\]

objects.
\end{theorem}

The notation

\[
\lceil x\rceil
\]

is read ``the ceiling of \(x\)'' and means the least integer greater
than or equal to \(x\).

Similarly,

\[
\lfloor x\rfloor
\]

is read ``the floor of \(x\)'' and means the greatest integer less than
or equal to \(x\).

%%%%%%%%%%%%%%%%%%%%%%%%%%%%%%%%%%%%%%%%%%%%%%%%%%%%%%%%%%%%
\subsection{Worked Example: Generalized Pigeonhole Principle}
\label{subsec:worked-generalized-pigeonhole}
%%%%%%%%%%%%%%%%%%%%%%%%%%%%%%%%%%%%%%%%%%%%%%%%%%%%%%%%%%%%

\begin{workedexamplebox}{Assignments Among Teams}

Suppose \(47\) tasks are assigned to \(6\) teams.

Then some team receives at least

\[
\left\lceil
\frac{47}{6}
\right\rceil
=
8
\]

tasks.

\begin{answerbox}
\[
\boxed{
\text{At least one team receives at least }8\text{ tasks.}
}
\]
\end{answerbox}

\end{workedexamplebox}

%%%%%%%%%%%%%%%%%%%%%%%%%%%%%%%%%%%%%%%%%%%%%%%%%%%%%%%%%%%%
\subsection{Factorials}
\label{subsec:factorials}
%%%%%%%%%%%%%%%%%%%%%%%%%%%%%%%%%%%%%%%%%%%%%%%%%%%%%%%%%%%%

For a positive integer \(n\), define

\[
\boxed{
n!
=
n(n-1)(n-2)\cdots2\cdot1.
}
\]

The symbol \(n!\) is read ``\(n\) factorial.''

By convention,

\[
\boxed{
0!=1.
}
\]

Examples include

\[
3!=6,
\]

\[
5!=120,
\]

and

\[
7!=5040.
\]

Factorials arise whenever distinct objects are placed in order.

%%%%%%%%%%%%%%%%%%%%%%%%%%%%%%%%%%%%%%%%%%%%%%%%%%%%%%%%%%%%
\subsection{Permutations}
\label{subsec:permutations-basic-counting}
%%%%%%%%%%%%%%%%%%%%%%%%%%%%%%%%%%%%%%%%%%%%%%%%%%%%%%%%%%%%

\begin{definitionbox}{Permutation}
A permutation of \(n\) distinct objects is an ordering of all \(n\)
objects.
\end{definitionbox}

There are

\[
n
\]

choices for the first position,

\[
n-1
\]

for the second,

and so on.

Therefore,

\[
\boxed{
n!
}
\]

permutations exist.

%%%%%%%%%%%%%%%%%%%%%%%%%%%%%%%%%%%%%%%%%%%%%%%%%%%%%%%%%%%%
\subsection{Worked Example: Permuting Objects}
\label{subsec:worked-permutations}
%%%%%%%%%%%%%%%%%%%%%%%%%%%%%%%%%%%%%%%%%%%%%%%%%%%%%%%%%%%%

\begin{workedexamplebox}{Arranging Five Books}

Five distinct books are placed in a row.

The first position has \(5\) choices.

The next has \(4\).

Continuing gives

\[
5\cdot4\cdot3\cdot2\cdot1
=
5!.
\]

\begin{answerbox}
\[
\boxed{
5!=120
}
\]

arrangements are possible.
\end{answerbox}

\end{workedexamplebox}

%%%%%%%%%%%%%%%%%%%%%%%%%%%%%%%%%%%%%%%%%%%%%%%%%%%%%%%%%%%%
\subsection{Ordered Selections}
\label{subsec:ordered-selections}
%%%%%%%%%%%%%%%%%%%%%%%%%%%%%%%%%%%%%%%%%%%%%%%%%%%%%%%%%%%%

Suppose \(r\) distinct objects are selected from \(n\) distinct objects
and order matters.

The number of choices is

\[
n(n-1)\cdots(n-r+1).
\]

This may be written

\[
\boxed{
P(n,r)
=
\frac{n!}{(n-r)!}.
}
\]

Other common notations include

\[
{}_nP_r
\]

and

\[
(n)_r.
\]

The expression

\[
(n)_r
\]

is often called a falling factorial.

%%%%%%%%%%%%%%%%%%%%%%%%%%%%%%%%%%%%%%%%%%%%%%%%%%%%%%%%%%%%
\subsection{Falling Factorials}
\label{subsec:falling-factorials}
%%%%%%%%%%%%%%%%%%%%%%%%%%%%%%%%%%%%%%%%%%%%%%%%%%%%%%%%%%%%

Define

\[
\boxed{
(n)_r
=
n(n-1)(n-2)\cdots(n-r+1).
}
\]

Thus,

\[
\boxed{
(n)_r
=
\frac{n!}{(n-r)!}.
}
\]

The falling factorial is useful because it directly expresses the
number of injective choices from \(r\) ordered positions into \(n\)
available objects.

%%%%%%%%%%%%%%%%%%%%%%%%%%%%%%%%%%%%%%%%%%%%%%%%%%%%%%%%%%%%
\subsection{Worked Example: Ordered Selection}
\label{subsec:worked-ordered-selection}
%%%%%%%%%%%%%%%%%%%%%%%%%%%%%%%%%%%%%%%%%%%%%%%%%%%%%%%%%%%%

\begin{workedexamplebox}{Awarding Three Places}

Ten contestants compete for first, second, and third place.

The positions are distinct, so order matters.

The number of results is

\[
10\cdot9\cdot8.
\]

Equivalently,

\[
P(10,3)
=
\frac{10!}{7!}.
\]

\begin{answerbox}
\[
\boxed{
720
}
\]

possible ordered results exist.
\end{answerbox}

\end{workedexamplebox}

%%%%%%%%%%%%%%%%%%%%%%%%%%%%%%%%%%%%%%%%%%%%%%%%%%%%%%%%%%%%
\subsection{Combinations}
\label{subsec:combinations}
%%%%%%%%%%%%%%%%%%%%%%%%%%%%%%%%%%%%%%%%%%%%%%%%%%%%%%%%%%%%

Suppose \(r\) objects are selected from \(n\) distinct objects and order
does not matter.

\begin{definitionbox}{Binomial Coefficient}
The number of \(r\)-element subsets of an \(n\)-element set is

\[
\boxed{
\binom{n}{r}
=
\frac{n!}{r!(n-r)!}.
}
\]
\end{definitionbox}

The expression

\[
\binom{n}{r}
\]

is read ``\(n\) choose \(r\).''

It is also called a \emph{binomial coefficient}.

%%%%%%%%%%%%%%%%%%%%%%%%%%%%%%%%%%%%%%%%%%%%%%%%%%%%%%%%%%%%
\subsection{Why the Combination Formula Works}
\label{subsec:why-combination-formula}
%%%%%%%%%%%%%%%%%%%%%%%%%%%%%%%%%%%%%%%%%%%%%%%%%%%%%%%%%%%%

There are

\[
\frac{n!}{(n-r)!}
\]

ordered selections of \(r\) objects.

Each unordered \(r\)-element subset can be ordered in

\[
r!
\]

ways.

Therefore, by the division principle,

\[
\binom{n}{r}
=
\frac{1}{r!}
\frac{n!}{(n-r)!}.
\]

Hence,

\[
\boxed{
\binom{n}{r}
=
\frac{n!}{r!(n-r)!}.
}
\]

%%%%%%%%%%%%%%%%%%%%%%%%%%%%%%%%%%%%%%%%%%%%%%%%%%%%%%%%%%%%
\subsection{Worked Example: Committee Selection}
\label{subsec:worked-committee-selection}
%%%%%%%%%%%%%%%%%%%%%%%%%%%%%%%%%%%%%%%%%%%%%%%%%%%%%%%%%%%%

\begin{workedexamplebox}{Choosing a Committee}

A committee of \(4\) people is selected from \(10\) people.

The order in which committee members are selected is irrelevant.

Thus,

\[
\binom{10}{4}
=
\frac{10!}{4!6!}.
\]

\begin{answerbox}
\[
\boxed{
\binom{10}{4}=210
}
\]

committees are possible.
\end{answerbox}

\end{workedexamplebox}

%%%%%%%%%%%%%%%%%%%%%%%%%%%%%%%%%%%%%%%%%%%%%%%%%%%%%%%%%%%%
\subsection{Symmetry of Binomial Coefficients}
\label{subsec:binomial-symmetry}
%%%%%%%%%%%%%%%%%%%%%%%%%%%%%%%%%%%%%%%%%%%%%%%%%%%%%%%%%%%%

Choosing \(r\) elements to include is equivalent to choosing the
remaining \(n-r\) elements to exclude.

Therefore,

\[
\boxed{
\binom{n}{r}
=
\binom{n}{n-r}.
}
\]

Algebraically,

\[
\frac{n!}{r!(n-r)!}
=
\frac{n!}{(n-r)!r!}.
\]

This is one of the simplest examples of a combinatorial identity with a
direct bijective explanation.

%%%%%%%%%%%%%%%%%%%%%%%%%%%%%%%%%%%%%%%%%%%%%%%%%%%%%%%%%%%%
\subsection{Boundary Values}
\label{subsec:binomial-boundary-values}
%%%%%%%%%%%%%%%%%%%%%%%%%%%%%%%%%%%%%%%%%%%%%%%%%%%%%%%%%%%%

For every nonnegative integer \(n\),

\[
\boxed{
\binom{n}{0}=1
}
\]

and

\[
\boxed{
\binom{n}{n}=1.
}
\]

There is exactly one way to choose no elements and exactly one way to
choose all elements.

Also,

\[
\boxed{
\binom{n}{1}=n.
}
\]

%%%%%%%%%%%%%%%%%%%%%%%%%%%%%%%%%%%%%%%%%%%%%%%%%%%%%%%%%%%%
\subsection{Pascal's Identity}
\label{subsec:pascals-identity}
%%%%%%%%%%%%%%%%%%%%%%%%%%%%%%%%%%%%%%%%%%%%%%%%%%%%%%%%%%%%

\begin{theorem}[Pascal's Identity]
For suitable integers \(n\) and \(r\),

\[
\boxed{
\binom{n}{r}
=
\binom{n-1}{r}
+
\binom{n-1}{r-1}.
}
\]
\end{theorem}

To see why, fix one distinguished element \(x\).

Every \(r\)-element subset either:

\begin{itemize}
    \item excludes \(x\), giving
    \[
    \binom{n-1}{r}
    \]
    possibilities; or

    \item includes \(x\), requiring
    \[
    r-1
    \]
    additional elements from the remaining \(n-1\), giving
    \[
    \binom{n-1}{r-1}.
    \]
\end{itemize}

The two cases are disjoint.

%%%%%%%%%%%%%%%%%%%%%%%%%%%%%%%%%%%%%%%%%%%%%%%%%%%%%%%%%%%%
\subsection{Pascal's Triangle}
\label{subsec:pascals-triangle}
%%%%%%%%%%%%%%%%%%%%%%%%%%%%%%%%%%%%%%%%%%%%%%%%%%%%%%%%%%%%

Pascal's identity generates the familiar triangular array

\[
\begin{array}{ccccccccc}
&&&&1&&&&\\
&&&1&&1&&&\\
&&1&&2&&1&&\\
&1&&3&&3&&1&\\
1&&4&&6&&4&&1
\end{array}
\]

Each interior entry equals the sum of the two entries above it.

The \(n\)-th row contains the coefficients

\[
\boxed{
\binom{n}{0},
\binom{n}{1},
\ldots,
\binom{n}{n}.
}
\]

%%%%%%%%%%%%%%%%%%%%%%%%%%%%%%%%%%%%%%%%%%%%%%%%%%%%%%%%%%%%
\subsection{The Binomial Theorem}
\label{subsec:binomial-theorem-basic}
%%%%%%%%%%%%%%%%%%%%%%%%%%%%%%%%%%%%%%%%%%%%%%%%%%%%%%%%%%%%

The binomial coefficients also arise algebraically.

\begin{theorem}[Binomial Theorem]
For every nonnegative integer \(n\),

\[
\boxed{
(x+y)^n
=
\sum_{k=0}^{n}
\binom{n}{k}
x^{n-k}y^k.
}
\]
\end{theorem}

The coefficient

\[
\binom{n}{k}
\]

appears because to produce the term

\[
x^{n-k}y^k
\]

from

\[
(x+y)^n,
\]

one must choose the \(y\)-term from exactly \(k\) of the \(n\)
factors.

%%%%%%%%%%%%%%%%%%%%%%%%%%%%%%%%%%%%%%%%%%%%%%%%%%%%%%%%%%%%
\subsection{Worked Example: Binomial Expansion}
\label{subsec:worked-binomial-expansion}
%%%%%%%%%%%%%%%%%%%%%%%%%%%%%%%%%%%%%%%%%%%%%%%%%%%%%%%%%%%%

\begin{workedexamplebox}{Expanding a Fourth Power}

Using the binomial theorem,

\[
(x+y)^4
=
\binom{4}{0}x^4
+
\binom{4}{1}x^3y
+
\binom{4}{2}x^2y^2
+
\binom{4}{3}xy^3
+
\binom{4}{4}y^4.
\]

Therefore,

\begin{answerbox}
\[
\boxed{
(x+y)^4
=
x^4+4x^3y+6x^2y^2+4xy^3+y^4.
}
\]
\end{answerbox}

\end{workedexamplebox}

%%%%%%%%%%%%%%%%%%%%%%%%%%%%%%%%%%%%%%%%%%%%%%%%%%%%%%%%%%%%
\subsection{Sum of Binomial Coefficients}
\label{subsec:sum-binomial-coefficients}
%%%%%%%%%%%%%%%%%%%%%%%%%%%%%%%%%%%%%%%%%%%%%%%%%%%%%%%%%%%%

Set

\[
x=y=1
\]

in the binomial theorem.

Then

\[
2^n
=
\sum_{k=0}^{n}\binom{n}{k}.
\]

Thus,

\[
\boxed{
\sum_{k=0}^{n}\binom{n}{k}
=
2^n.
}
\]

Combinatorially, this counts all subsets of an \(n\)-element set by
grouping them according to their size.

%%%%%%%%%%%%%%%%%%%%%%%%%%%%%%%%%%%%%%%%%%%%%%%%%%%%%%%%%%%%
\subsection{Alternating Sum of Binomial Coefficients}
\label{subsec:alternating-binomial-sum}
%%%%%%%%%%%%%%%%%%%%%%%%%%%%%%%%%%%%%%%%%%%%%%%%%%%%%%%%%%%%

Set

\[
x=1,
\qquad
y=-1
\]

in the binomial theorem.

For \(n\geq1\),

\[
0
=
\sum_{k=0}^{n}
(-1)^k\binom{n}{k}.
\]

Hence,

\[
\boxed{
\sum_{k=0}^{n}
(-1)^k\binom{n}{k}
=
0.
}
\]

This identity will later appear naturally in inclusion--exclusion and
Möbius inversion.

%%%%%%%%%%%%%%%%%%%%%%%%%%%%%%%%%%%%%%%%%%%%%%%%%%%%%%%%%%%%
\subsection{Multinomial Coefficients}
\label{subsec:multinomial-coefficients}
%%%%%%%%%%%%%%%%%%%%%%%%%%%%%%%%%%%%%%%%%%%%%%%%%%%%%%%%%%%%

The binomial coefficient extends naturally to more than two groups.

Suppose

\[
n_1+n_2+\cdots+n_k=n.
\]

The multinomial coefficient is

\[
\boxed{
\binom{n}{n_1,n_2,\ldots,n_k}
=
\frac{n!}{n_1!n_2!\cdots n_k!}.
}
\]

It counts the number of ways to divide \(n\) distinct objects into
ordered groups of sizes

\[
n_1,\ldots,n_k.
\]

%%%%%%%%%%%%%%%%%%%%%%%%%%%%%%%%%%%%%%%%%%%%%%%%%%%%%%%%%%%%
\subsection{Permutations with Repeated Objects}
\label{subsec:permutations-repeated-objects}
%%%%%%%%%%%%%%%%%%%%%%%%%%%%%%%%%%%%%%%%%%%%%%%%%%%%%%%%%%%%

Suppose \(n\) objects contain repeated types.

If there are:

\[
n_1
\]

identical objects of the first type,

\[
n_2
\]

of the second type,

and so on, with

\[
n_1+\cdots+n_k=n,
\]

then the number of distinct permutations is

\[
\boxed{
\frac{n!}{n_1!n_2!\cdots n_k!}.
}
\]

%%%%%%%%%%%%%%%%%%%%%%%%%%%%%%%%%%%%%%%%%%%%%%%%%%%%%%%%%%%%
\subsection{Worked Example: Repeated Letters}
\label{subsec:worked-repeated-letters}
%%%%%%%%%%%%%%%%%%%%%%%%%%%%%%%%%%%%%%%%%%%%%%%%%%%%%%%%%%%%

\begin{workedexamplebox}{Arrangements of MATHEMATICS}

Consider the word

\[
\text{MATHEMATICS}.
\]

It contains \(11\) letters.

The repeated letters are:

\[
M:2,
\qquad
A:2,
\qquad
T:2.
\]

The remaining letters occur once.

Therefore,

\[
\frac{11!}{2!2!2!}.
\]

\begin{answerbox}
\[
\boxed{
\frac{11!}{2!2!2!}
=
4,989,600
}
\]

distinct arrangements are possible.
\end{answerbox}

\end{workedexamplebox}

%%%%%%%%%%%%%%%%%%%%%%%%%%%%%%%%%%%%%%%%%%%%%%%%%%%%%%%%%%%%
\subsection{Ordered Sampling with Repetition}
\label{subsec:ordered-sampling-repetition}
%%%%%%%%%%%%%%%%%%%%%%%%%%%%%%%%%%%%%%%%%%%%%%%%%%%%%%%%%%%%

Suppose \(r\) selections are made from \(n\) available types.

If repetition is allowed and order matters, each position has \(n\)
choices.

Therefore,

\[
\boxed{
n^r
}
\]

ordered selections are possible.

This is the same structure as strings of length \(r\) over an
\(n\)-symbol alphabet.

%%%%%%%%%%%%%%%%%%%%%%%%%%%%%%%%%%%%%%%%%%%%%%%%%%%%%%%%%%%%
\subsection{Selection Summary}
\label{subsec:selection-summary}
%%%%%%%%%%%%%%%%%%%%%%%%%%%%%%%%%%%%%%%%%%%%%%%%%%%%%%%%%%%%

For \(r\) selections from \(n\) distinct types:

\[
\boxed{
\begin{array}{c|c|c}
&\text{Repetition Not Allowed}
&\text{Repetition Allowed}\\
\hline
\text{Order Matters}
&
\dfrac{n!}{(n-r)!}
&
n^r
\\[3mm]
\text{Order Does Not Matter}
&
\binom{n}{r}
&
\binom{n+r-1}{r}
\end{array}
}
\]

The final entry will be developed using stars and bars.

%%%%%%%%%%%%%%%%%%%%%%%%%%%%%%%%%%%%%%%%%%%%%%%%%%%%%%%%%%%%
\subsection{Integer Solutions and Stars and Bars}
\label{subsec:stars-and-bars}
%%%%%%%%%%%%%%%%%%%%%%%%%%%%%%%%%%%%%%%%%%%%%%%%%%%%%%%%%%%%

Suppose we seek nonnegative integer solutions of

\[
x_1+x_2+\cdots+x_k=n.
\]

Imagine \(n\) identical objects represented by stars.

To divide them into \(k\) groups, insert

\[
k-1
\]

bars.

For example,

\[
**|****|*
\]

represents a distribution among three groups.

The total number of symbols is

\[
n+k-1.
\]

We choose which \(k-1\) positions contain bars.

Therefore,

\[
\boxed{
\#\{(x_1,\ldots,x_k)\in\mathbb{Z}_{\geq0}^k:
x_1+\cdots+x_k=n\}
=
\binom{n+k-1}{k-1}.
}
\]

Equivalently,

\[
\boxed{
\binom{n+k-1}{n}.
}
\]

%%%%%%%%%%%%%%%%%%%%%%%%%%%%%%%%%%%%%%%%%%%%%%%%%%%%%%%%%%%%
\subsection{Worked Example: Stars and Bars}
\label{subsec:worked-stars-and-bars}
%%%%%%%%%%%%%%%%%%%%%%%%%%%%%%%%%%%%%%%%%%%%%%%%%%%%%%%%%%%%

\begin{workedexamplebox}{Distributing Identical Objects}

How many ways can \(10\) identical objects be distributed among \(4\)
distinct boxes if empty boxes are allowed?

We seek nonnegative integer solutions to

\[
x_1+x_2+x_3+x_4=10.
\]

Thus,

\[
\binom{10+4-1}{4-1}
=
\binom{13}{3}.
\]

\begin{answerbox}
\[
\boxed{
\binom{13}{3}=286
}
\]

distributions are possible.
\end{answerbox}

\end{workedexamplebox}

%%%%%%%%%%%%%%%%%%%%%%%%%%%%%%%%%%%%%%%%%%%%%%%%%%%%%%%%%%%%
\subsection{Positive Integer Solutions}
\label{subsec:positive-integer-solutions}
%%%%%%%%%%%%%%%%%%%%%%%%%%%%%%%%%%%%%%%%%%%%%%%%%%%%%%%%%%%%

Suppose instead that

\[
x_i\geq1
\]

for every \(i\).

Set

\[
y_i=x_i-1.
\]

Then

\[
y_i\geq0
\]

and

\[
y_1+\cdots+y_k=n-k.
\]

Therefore,

\[
\boxed{
\#\{x_1+\cdots+x_k=n,\ x_i\geq1\}
=
\binom{n-1}{k-1}.
}
\]

This formula assumes \(n\geq k\).

%%%%%%%%%%%%%%%%%%%%%%%%%%%%%%%%%%%%%%%%%%%%%%%%%%%%%%%%%%%%
\subsection{Worked Example: Positive Distributions}
\label{subsec:worked-positive-distributions}
%%%%%%%%%%%%%%%%%%%%%%%%%%%%%%%%%%%%%%%%%%%%%%%%%%%%%%%%%%%%

\begin{workedexamplebox}{At Least One Object per Box}

How many ways can \(12\) identical objects be distributed among \(4\)
distinct boxes if every box must receive at least one object?

We seek positive integer solutions of

\[
x_1+x_2+x_3+x_4=12.
\]

Therefore,

\[
\binom{12-1}{4-1}
=
\binom{11}{3}.
\]

\begin{answerbox}
\[
\boxed{
165
}
\]

distributions are possible.
\end{answerbox}

\end{workedexamplebox}

%%%%%%%%%%%%%%%%%%%%%%%%%%%%%%%%%%%%%%%%%%%%%%%%%%%%%%%%%%%%
\subsection{Lower Bounds in Stars and Bars}
\label{subsec:stars-bars-lower-bounds}
%%%%%%%%%%%%%%%%%%%%%%%%%%%%%%%%%%%%%%%%%%%%%%%%%%%%%%%%%%%%

Suppose

\[
x_i\geq a_i.
\]

Introduce

\[
y_i=x_i-a_i.
\]

Then

\[
y_i\geq0.
\]

The equation

\[
x_1+\cdots+x_k=n
\]

becomes

\[
y_1+\cdots+y_k
=
n-\sum_{i=1}^{k}a_i.
\]

Hence the lower bounds can be removed by a change of variables.

%%%%%%%%%%%%%%%%%%%%%%%%%%%%%%%%%%%%%%%%%%%%%%%%%%%%%%%%%%%%
\subsection{Combinations with Repetition}
\label{subsec:combinations-with-repetition}
%%%%%%%%%%%%%%%%%%%%%%%%%%%%%%%%%%%%%%%%%%%%%%%%%%%%%%%%%%%%

Suppose \(r\) objects are selected from \(n\) types, repetition is
allowed, and order does not matter.

Let

\[
x_i
\]

be the number selected of type \(i\).

Then

\[
x_1+\cdots+x_n=r
\]

with

\[
x_i\geq0.
\]

Therefore,

\[
\boxed{
\binom{n+r-1}{r}
}
\]

such selections exist.

This is sometimes called the number of \(r\)-combinations with
repetition from \(n\) types.

%%%%%%%%%%%%%%%%%%%%%%%%%%%%%%%%%%%%%%%%%%%%%%%%%%%%%%%%%%%%
\subsection{Upper Bounds and Why Stars and Bars May Fail}
\label{subsec:stars-bars-upper-bounds}
%%%%%%%%%%%%%%%%%%%%%%%%%%%%%%%%%%%%%%%%%%%%%%%%%%%%%%%%%%%%

Stars and bars handles lower bounds naturally.

Upper bounds require more care.

For example,

\[
x_1+x_2+x_3=12
\]

with

\[
0\leq x_i\leq5
\]

cannot be solved by the basic stars-and-bars formula alone.

One must subtract solutions violating the upper bounds.

This leads naturally to inclusion--exclusion.

%%%%%%%%%%%%%%%%%%%%%%%%%%%%%%%%%%%%%%%%%%%%%%%%%%%%%%%%%%%%
\subsection{The Principle of Inclusion--Exclusion}
\label{subsec:inclusion-exclusion}
%%%%%%%%%%%%%%%%%%%%%%%%%%%%%%%%%%%%%%%%%%%%%%%%%%%%%%%%%%%%

When sets overlap, the addition principle overcounts their
intersections.

For two finite sets,

\[
\boxed{
|A\cup B|
=
|A|+|B|-|A\cap B|.
}
\]

For three sets,

\[
\boxed{
\begin{aligned}
|A\cup B\cup C|
={}&
|A|+|B|+|C|\\
&-|A\cap B|
-|A\cap C|
-|B\cap C|\\
&+|A\cap B\cap C|.
\end{aligned}
}
\]

This alternating process generalizes.

%%%%%%%%%%%%%%%%%%%%%%%%%%%%%%%%%%%%%%%%%%%%%%%%%%%%%%%%%%%%
\subsection{General Inclusion--Exclusion}
\label{subsec:general-inclusion-exclusion}
%%%%%%%%%%%%%%%%%%%%%%%%%%%%%%%%%%%%%%%%%%%%%%%%%%%%%%%%%%%%

\begin{theorem}[Principle of Inclusion--Exclusion]
For finite sets \(A_1,\ldots,A_n\),

\[
\boxed{
\left|
\bigcup_{i=1}^{n}A_i
\right|
=
\sum_i|A_i|
-
\sum_{i<j}|A_i\cap A_j|
+
\sum_{i<j<k}|A_i\cap A_j\cap A_k|
-\cdots
+
(-1)^{n+1}
\left|
\bigcap_{i=1}^{n}A_i
\right|.
}
\]
\end{theorem}

The principle is often abbreviated PIE.

%%%%%%%%%%%%%%%%%%%%%%%%%%%%%%%%%%%%%%%%%%%%%%%%%%%%%%%%%%%%
\subsection{Why Inclusion--Exclusion Works}
\label{subsec:why-inclusion-exclusion-works}
%%%%%%%%%%%%%%%%%%%%%%%%%%%%%%%%%%%%%%%%%%%%%%%%%%%%%%%%%%%%

Suppose an element belongs to exactly \(r\) of the sets.

Its total contribution to the right-hand side is

\[
\binom{r}{1}
-
\binom{r}{2}
+
\binom{r}{3}
-\cdots.
\]

Using

\[
\sum_{k=0}^{r}
(-1)^k\binom{r}{k}
=
0,
\]

we obtain

\[
\sum_{k=1}^{r}
(-1)^{k+1}\binom{r}{k}
=
1.
\]

Thus every element in the union is counted exactly once.

%%%%%%%%%%%%%%%%%%%%%%%%%%%%%%%%%%%%%%%%%%%%%%%%%%%%%%%%%%%%
\subsection{Worked Example: Divisibility Restrictions}
\label{subsec:worked-inclusion-exclusion}
%%%%%%%%%%%%%%%%%%%%%%%%%%%%%%%%%%%%%%%%%%%%%%%%%%%%%%%%%%%%

\begin{workedexamplebox}{Integers Divisible by \(2\) or \(3\)}

How many integers from \(1\) through \(100\) are divisible by \(2\) or
\(3\)?

Let

\[
A=\{n:2\mid n\}
\]

and

\[
B=\{n:3\mid n\}.
\]

Then

\[
|A|=50,
\]

\[
|B|=33,
\]

and

\[
|A\cap B|=16
\]

because the intersection consists of multiples of \(6\).

Therefore,

\[
|A\cup B|
=
50+33-16.
\]

\begin{answerbox}
\[
\boxed{
67
}
\]

integers are divisible by \(2\) or \(3\).
\end{answerbox}

\end{workedexamplebox}

%%%%%%%%%%%%%%%%%%%%%%%%%%%%%%%%%%%%%%%%%%%%%%%%%%%%%%%%%%%%
\subsection{Counting Objects Avoiding Properties}
\label{subsec:counting-avoiding-properties}
%%%%%%%%%%%%%%%%%%%%%%%%%%%%%%%%%%%%%%%%%%%%%%%%%%%%%%%%%%%%

Suppose \(U\) is the universe of all candidate objects.

Let

\[
A_i
\]

be the set of objects violating condition \(i\).

Then the number satisfying every condition is

\[
\boxed{
|U|
-
\left|
\bigcup_iA_i
\right|.
}
\]

Inclusion--exclusion therefore becomes a powerful method for counting
objects that avoid several forbidden properties.

%%%%%%%%%%%%%%%%%%%%%%%%%%%%%%%%%%%%%%%%%%%%%%%%%%%%%%%%%%%%
\subsection{Derangements}
\label{subsec:derangements}
%%%%%%%%%%%%%%%%%%%%%%%%%%%%%%%%%%%%%%%%%%%%%%%%%%%%%%%%%%%%

A derangement is a permutation having no fixed points.

\begin{definitionbox}{Derangement}
A derangement of \(n\) objects is a permutation

\[
\pi
\]

such that

\[
\boxed{
\pi(i)\neq i
}
\]

for every \(i\).
\end{definitionbox}

The number of derangements of \(n\) objects is commonly denoted

\[
D_n
\]

or

\[
!n.
\]

%%%%%%%%%%%%%%%%%%%%%%%%%%%%%%%%%%%%%%%%%%%%%%%%%%%%%%%%%%%%
\subsection{Counting Derangements}
\label{subsec:counting-derangements}
%%%%%%%%%%%%%%%%%%%%%%%%%%%%%%%%%%%%%%%%%%%%%%%%%%%%%%%%%%%%

Let \(A_i\) be the set of permutations fixing element \(i\).

Then

\[
|A_i|=(n-1)!.
\]

For distinct \(i,j\),

\[
|A_i\cap A_j|=(n-2)!.
\]

Continuing with inclusion--exclusion gives

\[
\boxed{
D_n
=
n!
\sum_{k=0}^{n}\frac{(-1)^k}{k!}.
}
\]

Equivalently,

\[
\boxed{
D_n
=
n!
\left(
1-\frac1{1!}
+\frac1{2!}
-\frac1{3!}
+\cdots
+\frac{(-1)^n}{n!}
\right).
}
\]

%%%%%%%%%%%%%%%%%%%%%%%%%%%%%%%%%%%%%%%%%%%%%%%%%%%%%%%%%%%%
\subsection{Approximation for Derangements}
\label{subsec:derangement-approximation}
%%%%%%%%%%%%%%%%%%%%%%%%%%%%%%%%%%%%%%%%%%%%%%%%%%%%%%%%%%%%

Since

\[
e^{-1}
=
\sum_{k=0}^{\infty}
\frac{(-1)^k}{k!},
\]

the number of derangements satisfies approximately

\[
\boxed{
D_n\approx\frac{n!}{e}.
}
\]

In fact, \(D_n\) is the nearest integer to

\[
\frac{n!}{e}.
\]

This is an early example of an exact combinatorial count connected to
analytic constants.

%%%%%%%%%%%%%%%%%%%%%%%%%%%%%%%%%%%%%%%%%%%%%%%%%%%%%%%%%%%%
\subsection{Worked Example: Derangements}
\label{subsec:worked-derangements}
%%%%%%%%%%%%%%%%%%%%%%%%%%%%%%%%%%%%%%%%%%%%%%%%%%%%%%%%%%%%

\begin{workedexamplebox}{Derangements of Four Objects}

Using the formula,

\[
D_4
=
4!
\left(
1-1+\frac12-\frac16+\frac1{24}
\right).
\]

Thus,

\[
D_4
=
24
\left(
\frac{9}{24}
\right).
\]

\begin{answerbox}
\[
\boxed{
D_4=9.
}
\]
\end{answerbox}

\end{workedexamplebox}

%%%%%%%%%%%%%%%%%%%%%%%%%%%%%%%%%%%%%%%%%%%%%%%%%%%%%%%%%%%%
\subsection{Counting Functions}
\label{subsec:counting-functions}
%%%%%%%%%%%%%%%%%%%%%%%%%%%%%%%%%%%%%%%%%%%%%%%%%%%%%%%%%%%%

Suppose

\[
|A|=m
\]

and

\[
|B|=n.
\]

The number of all functions

\[
f:A\to B
\]

is

\[
\boxed{
n^m.
}
\]

Each element of \(A\) independently chooses one of \(n\) outputs.

%%%%%%%%%%%%%%%%%%%%%%%%%%%%%%%%%%%%%%%%%%%%%%%%%%%%%%%%%%%%
\subsection{Counting Injective Functions}
\label{subsec:counting-injective-functions}
%%%%%%%%%%%%%%%%%%%%%%%%%%%%%%%%%%%%%%%%%%%%%%%%%%%%%%%%%%%%

If

\[
m\leq n,
\]

the number of injective functions

\[
f:A\to B
\]

is

\[
\boxed{
n(n-1)\cdots(n-m+1)
=
\frac{n!}{(n-m)!}.
}
\]

If

\[
m>n,
\]

there are no injective functions, by the pigeonhole principle.

%%%%%%%%%%%%%%%%%%%%%%%%%%%%%%%%%%%%%%%%%%%%%%%%%%%%%%%%%%%%
\subsection{Counting Bijective Functions}
\label{subsec:counting-bijective-functions}
%%%%%%%%%%%%%%%%%%%%%%%%%%%%%%%%%%%%%%%%%%%%%%%%%%%%%%%%%%%%

If

\[
|A|=|B|=n,
\]

a bijection from \(A\) to \(B\) is simply an ordering of the possible
images.

Therefore, the number of bijections is

\[
\boxed{
n!.
}
\]

In particular, the number of permutations of an \(n\)-element set is

\[
n!.
\]

%%%%%%%%%%%%%%%%%%%%%%%%%%%%%%%%%%%%%%%%%%%%%%%%%%%%%%%%%%%%
\subsection{Counting Surjective Functions}
\label{subsec:counting-surjective-functions}
%%%%%%%%%%%%%%%%%%%%%%%%%%%%%%%%%%%%%%%%%%%%%%%%%%%%%%%%%%%%

Suppose

\[
|A|=m,
\qquad
|B|=n.
\]

To count surjections

\[
f:A\to B,
\]

begin with all

\[
n^m
\]

functions and exclude those that omit one or more elements of \(B\).

By inclusion--exclusion,

\[
\boxed{
\#\{\text{surjections }A\to B\}
=
\sum_{k=0}^{n}
(-1)^k
\binom{n}{k}
(n-k)^m.
}
\]

If \(m<n\), this count is zero.

%%%%%%%%%%%%%%%%%%%%%%%%%%%%%%%%%%%%%%%%%%%%%%%%%%%%%%%%%%%%
\subsection{Worked Example: Surjections}
\label{subsec:worked-surjections}
%%%%%%%%%%%%%%%%%%%%%%%%%%%%%%%%%%%%%%%%%%%%%%%%%%%%%%%%%%%%

\begin{workedexamplebox}{Onto Functions from Four Elements to Three}

Let

\[
|A|=4,
\qquad
|B|=3.
\]

The number of all functions is

\[
3^4.
\]

Subtract those omitting one target value:

\[
\binom31 2^4.
\]

Add back those omitting two target values:

\[
\binom32 1^4.
\]

Therefore,

\[
3^4-3(2^4)+3(1^4).
\]

\begin{answerbox}
\[
\boxed{
36
}
\]

surjective functions exist.
\end{answerbox}

\end{workedexamplebox}

%%%%%%%%%%%%%%%%%%%%%%%%%%%%%%%%%%%%%%%%%%%%%%%%%%%%%%%%%%%%
\subsection{Set Partitions and Stirling Numbers}
\label{subsec:set-partitions-stirling-preview}
%%%%%%%%%%%%%%%%%%%%%%%%%%%%%%%%%%%%%%%%%%%%%%%%%%%%%%%%%%%%

A partition of an \(n\)-element set into \(k\) nonempty unlabeled
blocks is counted by a Stirling number of the second kind,

\[
\boxed{
\left\{
\begin{matrix}
n\\
k
\end{matrix}
\right\}.
}
\]

This notation is read ``Stirling \(n\) over \(k\), second kind.''

The relationship with surjections is

\[
\boxed{
\#\{\text{surjections from an }n\text{-element set onto a }
k\text{-element set}\}
=
k!
\left\{
\begin{matrix}
n\\
k
\end{matrix}
\right\}.
}
\]

The factor \(k!\) labels the blocks by target elements.

Stirling numbers will be developed more fully in later combinatorial
sections.

%%%%%%%%%%%%%%%%%%%%%%%%%%%%%%%%%%%%%%%%%%%%%%%%%%%%%%%%%%%%
\subsection{Circular Permutations}
\label{subsec:circular-permutations}
%%%%%%%%%%%%%%%%%%%%%%%%%%%%%%%%%%%%%%%%%%%%%%%%%%%%%%%%%%%%

Suppose \(n\) distinct objects are arranged around a circle.

If rotations are considered equivalent, fixing one object removes the
rotational redundancy.

The remaining \(n-1\) objects can be arranged in

\[
\boxed{
(n-1)!
}
\]

ways.

Thus the number of circular permutations is

\[
\boxed{
(n-1)!.
}
\]

%%%%%%%%%%%%%%%%%%%%%%%%%%%%%%%%%%%%%%%%%%%%%%%%%%%%%%%%%%%%
\subsection{Worked Example: Circular Seating}
\label{subsec:worked-circular-seating}
%%%%%%%%%%%%%%%%%%%%%%%%%%%%%%%%%%%%%%%%%%%%%%%%%%%%%%%%%%%%

\begin{workedexamplebox}{Seating Six People at a Round Table}

Six distinct people sit around a circular table.

Rotating everyone together does not produce a new seating arrangement.

Fix one person.

The remaining five people can be placed in

\[
5!
\]

orders.

\begin{answerbox}
\[
\boxed{
5!=120
}
\]

circular arrangements are possible.
\end{answerbox}

\end{workedexamplebox}

%%%%%%%%%%%%%%%%%%%%%%%%%%%%%%%%%%%%%%%%%%%%%%%%%%%%%%%%%%%%
\subsection{Reflection and Circular Arrangements}
\label{subsec:reflection-circular-arrangements}
%%%%%%%%%%%%%%%%%%%%%%%%%%%%%%%%%%%%%%%%%%%%%%%%%%%%%%%%%%%%

Whether reflections count as identical depends on the problem.

For a round table, clockwise and counterclockwise arrangements are
usually considered different unless an additional reflection symmetry
is explicitly imposed.

For necklaces or bracelets, rotations and possibly reflections may be
considered equivalent.

Such problems require a systematic treatment of symmetry and group
actions and will be revisited in algebraic combinatorics.

\begin{warningbox}
Do not automatically divide a circular count by \(2\).

Reflections are equivalent only when the problem explicitly says that
mirror images represent the same object.
\end{warningbox}

%%%%%%%%%%%%%%%%%%%%%%%%%%%%%%%%%%%%%%%%%%%%%%%%%%%%%%%%%%%%
\subsection{Counting Subsets with Restrictions}
\label{subsec:counting-subsets-restrictions}
%%%%%%%%%%%%%%%%%%%%%%%%%%%%%%%%%%%%%%%%%%%%%%%%%%%%%%%%%%%%

Suppose an \(r\)-element subset is chosen from \(n\) objects, with one
distinguished element \(a\).

The number of subsets containing \(a\) is

\[
\boxed{
\binom{n-1}{r-1}.
}
\]

The number excluding \(a\) is

\[
\boxed{
\binom{n-1}{r}.
}
\]

Adding the two cases gives Pascal's identity.

%%%%%%%%%%%%%%%%%%%%%%%%%%%%%%%%%%%%%%%%%%%%%%%%%%%%%%%%%%%%
\subsection{Required and Forbidden Elements}
\label{subsec:required-forbidden-elements}
%%%%%%%%%%%%%%%%%%%%%%%%%%%%%%%%%%%%%%%%%%%%%%%%%%%%%%%%%%%%

Suppose an \(r\)-element subset is chosen from an \(n\)-element set.

If \(s\) specified elements must be included, choose the remaining

\[
r-s
\]

elements from the remaining

\[
n-s.
\]

Thus the count is

\[
\boxed{
\binom{n-s}{r-s}.
}
\]

If \(t\) specified elements are forbidden, choose all \(r\) elements
from

\[
n-t
\]

available objects:

\[
\boxed{
\binom{n-t}{r}.
}
\]

%%%%%%%%%%%%%%%%%%%%%%%%%%%%%%%%%%%%%%%%%%%%%%%%%%%%%%%%%%%%
\subsection{Worked Example: Committee Restrictions}
\label{subsec:worked-committee-restrictions}
%%%%%%%%%%%%%%%%%%%%%%%%%%%%%%%%%%%%%%%%%%%%%%%%%%%%%%%%%%%%

\begin{workedexamplebox}{Committee with Required Members}

A \(5\)-person committee is chosen from \(12\) people.

Two particular people must be included.

Those two positions are already determined.

Choose the remaining \(3\) members from the other \(10\):

\[
\binom{10}{3}.
\]

\begin{answerbox}
\[
\boxed{
120
}
\]

committees satisfy the requirement.
\end{answerbox}

\end{workedexamplebox}

%%%%%%%%%%%%%%%%%%%%%%%%%%%%%%%%%%%%%%%%%%%%%%%%%%%%%%%%%%%%
\subsection{Exactly One of Several Distinguished Objects}
\label{subsec:exactly-one-distinguished}
%%%%%%%%%%%%%%%%%%%%%%%%%%%%%%%%%%%%%%%%%%%%%%%%%%%%%%%%%%%%

Suppose an \(r\)-element subset is chosen from \(n\) objects and there
are \(k\) distinguished objects.

To choose exactly one distinguished object:

\[
\binom{k}{1}
\]

chooses the distinguished member, while

\[
\binom{n-k}{r-1}
\]

chooses the remaining ordinary elements.

Therefore,

\[
\boxed{
k\binom{n-k}{r-1}.
}
\]

%%%%%%%%%%%%%%%%%%%%%%%%%%%%%%%%%%%%%%%%%%%%%%%%%%%%%%%%%%%%
\subsection{At Least and At Most Conditions}
\label{subsec:at-least-at-most-counting}
%%%%%%%%%%%%%%%%%%%%%%%%%%%%%%%%%%%%%%%%%%%%%%%%%%%%%%%%%%%%

Words such as

\[
\text{at least}
\]

and

\[
\text{at most}
\]

usually signal a sum of cases or a complement.

For example,

\[
\text{at least one}
\]

is often easiest to count as

\[
\boxed{
\text{all possibilities}
-
\text{none}.
}
\]

Likewise,

\[
\text{at most }r
\]

often produces a sum

\[
\boxed{
\sum_{k=0}^{r}
\text{number of outcomes with exactly }k.
}
\]

%%%%%%%%%%%%%%%%%%%%%%%%%%%%%%%%%%%%%%%%%%%%%%%%%%%%%%%%%%%%
\subsection{Counting by Cases}
\label{subsec:counting-by-cases}
%%%%%%%%%%%%%%%%%%%%%%%%%%%%%%%%%%%%%%%%%%%%%%%%%%%%%%%%%%%%

A problem may be divided according to a natural statistic.

If every object belongs to exactly one case,

\[
S=S_1\cup S_2\cup\cdots\cup S_k
\]

with pairwise disjoint \(S_i\), then

\[
\boxed{
|S|
=
\sum_{i=1}^{k}|S_i|.
}
\]

Useful case variables include:

\begin{itemize}
    \item number of selected objects of a certain type;
    \item position of a distinguished object;
    \item parity;
    \item largest or smallest element;
    \item number of repeated values;
    \item number of occupied boxes.
\end{itemize}

%%%%%%%%%%%%%%%%%%%%%%%%%%%%%%%%%%%%%%%%%%%%%%%%%%%%%%%%%%%%
\subsection{Counting by the First Distinguished Event}
\label{subsec:first-distinguished-event}
%%%%%%%%%%%%%%%%%%%%%%%%%%%%%%%%%%%%%%%%%%%%%%%%%%%%%%%%%%%%

A useful strategy partitions objects according to the location of the
first occurrence of a specified event.

For example, binary strings containing at least one \(1\) may be divided
according to the position of their first \(1\).

If the first \(1\) occurs at position \(k\), then:

\[
k-1
\]

leading positions are forced to be zero,

the \(k\)-th position is \(1\),

and the remaining

\[
n-k
\]

positions are arbitrary.

Thus the number of such strings is

\[
2^{n-k}.
\]

Summing over \(k\) gives

\[
\sum_{k=1}^{n}2^{n-k}
=
2^n-1.
\]

%%%%%%%%%%%%%%%%%%%%%%%%%%%%%%%%%%%%%%%%%%%%%%%%%%%%%%%%%%%%
\subsection{Bijective Proofs}
\label{subsec:bijective-proofs}
%%%%%%%%%%%%%%%%%%%%%%%%%%%%%%%%%%%%%%%%%%%%%%%%%%%%%%%%%%%%

A bijective proof establishes an identity by constructing a bijection
between two sets.

Suppose finite sets \(A\) and \(B\) satisfy

\[
A\leftrightarrow B
\]

through a bijection.

Then

\[
\boxed{
|A|=|B|.
}
\]

The equality is explained structurally rather than algebraically.

%%%%%%%%%%%%%%%%%%%%%%%%%%%%%%%%%%%%%%%%%%%%%%%%%%%%%%%%%%%%
\subsection{Worked Example: Binomial Symmetry by Bijection}
\label{subsec:worked-binomial-bijection}
%%%%%%%%%%%%%%%%%%%%%%%%%%%%%%%%%%%%%%%%%%%%%%%%%%%%%%%%%%%%

\begin{workedexamplebox}{Choosing a Subset or Its Complement}

Let \(S\) be an \(n\)-element set.

Map every \(r\)-element subset

\[
A\subseteq S
\]

to its complement

\[
S\setminus A.
\]

The complement contains

\[
n-r
\]

elements.

The map is reversible.

Therefore it is a bijection between \(r\)-element subsets and
\((n-r)\)-element subsets.

\begin{answerbox}
\[
\boxed{
\binom{n}{r}
=
\binom{n}{n-r}.
}
\]
\end{answerbox}

\end{workedexamplebox}

%%%%%%%%%%%%%%%%%%%%%%%%%%%%%%%%%%%%%%%%%%%%%%%%%%%%%%%%%%%%
\subsection{Double Counting}
\label{subsec:double-counting}
%%%%%%%%%%%%%%%%%%%%%%%%%%%%%%%%%%%%%%%%%%%%%%%%%%%%%%%%%%%%

Another powerful method counts the same finite set in two different
ways.

Suppose \(S\) is counted using two different descriptions.

If one count gives \(A\) and another gives \(B\), then

\[
\boxed{
A=B.
}
\]

This technique is called double counting.

%%%%%%%%%%%%%%%%%%%%%%%%%%%%%%%%%%%%%%%%%%%%%%%%%%%%%%%%%%%%
\subsection{Worked Example: Handshaking Lemma}
\label{subsec:worked-double-counting-handshake}
%%%%%%%%%%%%%%%%%%%%%%%%%%%%%%%%%%%%%%%%%%%%%%%%%%%%%%%%%%%%

\begin{workedexamplebox}{Counting Vertex--Edge Incidences}

Let

\[
G=(V,E)
\]

be a finite undirected graph.

Count pairs

\[
(v,e)
\]

such that vertex \(v\) is incident with edge \(e\).

Counting by vertices gives

\[
\sum_{v\in V}\deg(v).
\]

Counting by edges gives

\[
2|E|,
\]

because every edge has two endpoints.

Therefore,

\begin{answerbox}
\[
\boxed{
\sum_{v\in V}\deg(v)=2|E|.
}
\]
\end{answerbox}

\end{workedexamplebox}

%%%%%%%%%%%%%%%%%%%%%%%%%%%%%%%%%%%%%%%%%%%%%%%%%%%%%%%%%%%%
\subsection{Vandermonde's Identity}
\label{subsec:vandermonde-identity}
%%%%%%%%%%%%%%%%%%%%%%%%%%%%%%%%%%%%%%%%%%%%%%%%%%%%%%%%%%%%

A classical identity is

\[
\boxed{
\binom{m+n}{r}
=
\sum_{k}
\binom{m}{k}
\binom{n}{r-k}.
}
\]

To interpret it, divide a set of \(m+n\) objects into two groups of
sizes \(m\) and \(n\).

An \(r\)-element subset may contain exactly \(k\) objects from the first
group and

\[
r-k
\]

from the second.

Summing over all possible \(k\) gives the identity.

%%%%%%%%%%%%%%%%%%%%%%%%%%%%%%%%%%%%%%%%%%%%%%%%%%%%%%%%%%%%
\subsection{Weighted Double Counting}
\label{subsec:weighted-double-counting}
%%%%%%%%%%%%%%%%%%%%%%%%%%%%%%%%%%%%%%%%%%%%%%%%%%%%%%%%%%%%

Double counting can be extended beyond ordinary cardinalities.

Instead of assigning weight \(1\) to each object, one may assign a
weight

\[
w(x)
\]

and compute

\[
\sum_{x\in S}w(x)
\]

in two different ways.

This technique appears frequently in generating functions,
probability, graph theory, and algebraic combinatorics.

%%%%%%%%%%%%%%%%%%%%%%%%%%%%%%%%%%%%%%%%%%%%%%%%%%%%%%%%%%%%
\subsection{Counting Lattice Paths}
\label{subsec:lattice-paths-basic}
%%%%%%%%%%%%%%%%%%%%%%%%%%%%%%%%%%%%%%%%%%%%%%%%%%%%%%%%%%%%

Consider paths from

\[
(0,0)
\]

to

\[
(m,n)
\]

using only:

\[
R=(1,0)
\]

and

\[
U=(0,1).
\]

Every path contains exactly

\[
m
\]

right steps and

\[
n
\]

up steps.

Thus every path is determined by choosing which \(n\) positions among
the

\[
m+n
\]

steps are upward moves.

Therefore,

\[
\boxed{
\binom{m+n}{n}
=
\binom{m+n}{m}
}
\]

such lattice paths exist.

%%%%%%%%%%%%%%%%%%%%%%%%%%%%%%%%%%%%%%%%%%%%%%%%%%%%%%%%%%%%
\subsection{Worked Example: Lattice Paths}
\label{subsec:worked-lattice-path}
%%%%%%%%%%%%%%%%%%%%%%%%%%%%%%%%%%%%%%%%%%%%%%%%%%%%%%%%%%%%

\begin{workedexamplebox}{Paths to \((4,3)\)}

A path from

\[
(0,0)
\]

to

\[
(4,3)
\]

requires \(4\) right steps and \(3\) up steps.

There are \(7\) total steps.

Choose the \(3\) positions occupied by upward steps:

\[
\binom73.
\]

\begin{answerbox}
\[
\boxed{
35
}
\]

such paths exist.
\end{answerbox}

\end{workedexamplebox}

%%%%%%%%%%%%%%%%%%%%%%%%%%%%%%%%%%%%%%%%%%%%%%%%%%%%%%%%%%%%
\subsection{Counting Binary Strings by Weight}
\label{subsec:binary-strings-by-weight}
%%%%%%%%%%%%%%%%%%%%%%%%%%%%%%%%%%%%%%%%%%%%%%%%%%%%%%%%%%%%

A binary string of length \(n\) having exactly \(k\) ones is determined
by choosing which \(k\) positions contain \(1\).

Therefore,

\[
\boxed{
\#\{x\in\{0,1\}^n:\operatorname{wt}(x)=k\}
=
\binom{n}{k}.
}
\]

This gives another interpretation of binomial coefficients.

%%%%%%%%%%%%%%%%%%%%%%%%%%%%%%%%%%%%%%%%%%%%%%%%%%%%%%%%%%%%
\subsection{The Boolean Cube by Layers}
\label{subsec:boolean-cube-layers}
%%%%%%%%%%%%%%%%%%%%%%%%%%%%%%%%%%%%%%%%%%%%%%%%%%%%%%%%%%%%

The discrete cube

\[
\{0,1\}^n
\]

can be partitioned into layers according to Hamming weight.

Define

\[
L_k
=
\{x\in\{0,1\}^n:\operatorname{wt}(x)=k\}.
\]

Then

\[
\boxed{
|L_k|=\binom{n}{k}.
}
\]

Since the layers partition the cube,

\[
\sum_{k=0}^{n}\binom{n}{k}
=
2^n.
\]

This structure becomes fundamental in extremal combinatorics.

%%%%%%%%%%%%%%%%%%%%%%%%%%%%%%%%%%%%%%%%%%%%%%%%%%%%%%%%%%%%
\subsection{Counting Paths Through an Intermediate Point}
\label{subsec:lattice-path-intermediate}
%%%%%%%%%%%%%%%%%%%%%%%%%%%%%%%%%%%%%%%%%%%%%%%%%%%%%%%%%%%%

Suppose a monotone lattice path from

\[
(0,0)
\]

to

\[
(m,n)
\]

must pass through

\[
(a,b).
\]

The number of paths from the start to \((a,b)\) is

\[
\binom{a+b}{a}.
\]

The number from \((a,b)\) to \((m,n)\) is

\[
\binom{(m-a)+(n-b)}{m-a}.
\]

By the multiplication principle, the total is

\[
\boxed{
\binom{a+b}{a}
\binom{m+n-a-b}{m-a}.
}
\]

%%%%%%%%%%%%%%%%%%%%%%%%%%%%%%%%%%%%%%%%%%%%%%%%%%%%%%%%%%%%
\subsection{The Multinomial Theorem}
\label{subsec:multinomial-theorem}
%%%%%%%%%%%%%%%%%%%%%%%%%%%%%%%%%%%%%%%%%%%%%%%%%%%%%%%%%%%%

The binomial theorem extends to several variables.

\begin{theorem}[Multinomial Theorem]
For nonnegative integer \(n\),

\[
\boxed{
(x_1+\cdots+x_k)^n
=
\sum_{n_1+\cdots+n_k=n}
\frac{n!}{n_1!\cdots n_k!}
x_1^{n_1}\cdots x_k^{n_k}.
}
\]
\end{theorem}

The sum is taken over all nonnegative integer solutions of

\[
n_1+\cdots+n_k=n.
\]

The coefficient

\[
\frac{n!}{n_1!\cdots n_k!}
\]

counts the ways to choose which factors contribute each variable.

%%%%%%%%%%%%%%%%%%%%%%%%%%%%%%%%%%%%%%%%%%%%%%%%%%%%%%%%%%%%
\subsection{Counting by Encoding}
\label{subsec:counting-by-encoding}
%%%%%%%%%%%%%%%%%%%%%%%%%%%%%%%%%%%%%%%%%%%%%%%%%%%%%%%%%%%%

A useful counting strategy is to encode each desired object as a simpler
object whose count is already known.

Examples include:

\[
\text{subset}
\longleftrightarrow
\text{binary string},
\]

\[
\text{lattice path}
\longleftrightarrow
\text{word in }R,U,
\]

and

\[
\text{distribution}
\longleftrightarrow
\text{stars-and-bars string}.
\]

If the encoding is bijective, counting either representation gives the
same answer.

%%%%%%%%%%%%%%%%%%%%%%%%%%%%%%%%%%%%%%%%%%%%%%%%%%%%%%%%%%%%
\subsection{Overcounting}
\label{subsec:overcounting}
%%%%%%%%%%%%%%%%%%%%%%%%%%%%%%%%%%%%%%%%%%%%%%%%%%%%%%%%%%%%

A common error occurs when one object has several representations.

For example, if one counts unordered pairs by choosing a first and
second object, then each pair is counted twice.

This is overcounting.

The standard correction is:

\[
\boxed{
\text{desired count}
=
\frac{\text{number of representations}}
{\text{representations per object}}.
}
\]

This method is valid only when every object has the same number of
representations.

%%%%%%%%%%%%%%%%%%%%%%%%%%%%%%%%%%%%%%%%%%%%%%%%%%%%%%%%%%%%
\subsection{Unequal Overcounting}
\label{subsec:unequal-overcounting}
%%%%%%%%%%%%%%%%%%%%%%%%%%%%%%%%%%%%%%%%%%%%%%%%%%%%%%%%%%%%

If different objects are represented different numbers of times, a
single division factor does not work.

This frequently occurs in symmetry problems.

For example, highly symmetric configurations may have more
automorphisms than generic configurations.

Such situations require more advanced tools involving group actions,
orbits, stabilizers, Burnside-type counting, or cycle indices.

These ideas belong naturally to algebraic combinatorics.

%%%%%%%%%%%%%%%%%%%%%%%%%%%%%%%%%%%%%%%%%%%%%%%%%%%%%%%%%%%%
\subsection{Common Counting Strategy}
\label{subsec:common-counting-strategy}
%%%%%%%%%%%%%%%%%%%%%%%%%%%%%%%%%%%%%%%%%%%%%%%%%%%%%%%%%%%%

A reliable counting workflow is:

\begin{enumerate}
    \item identify the exact object being counted;
    \item determine whether order matters;
    \item determine whether repetition is allowed;
    \item identify restrictions;
    \item decide whether to count directly or by complement;
    \item break the problem into disjoint cases if necessary;
    \item use multiplication for successive independent choices;
    \item correct any overcounting;
    \item verify that every valid object is counted exactly once.
\end{enumerate}

%%%%%%%%%%%%%%%%%%%%%%%%%%%%%%%%%%%%%%%%%%%%%%%%%%%%%%%%%%%%
\subsection{A Counting Decision Table}
\label{subsec:counting-decision-table}
%%%%%%%%%%%%%%%%%%%%%%%%%%%%%%%%%%%%%%%%%%%%%%%%%%%%%%%%%%%%

For selections of size \(r\) from \(n\) types:

\[
\boxed{
\begin{array}{c|c|c}
& \text{Without Repetition} & \text{With Repetition}\\
\hline
\text{Ordered}
&
\dfrac{n!}{(n-r)!}
&
n^r
\\[3mm]
\text{Unordered}
&
\binom{n}{r}
&
\binom{n+r-1}{r}
\end{array}
}
\]

This table is useful, but it should not replace understanding the
underlying structure.

%%%%%%%%%%%%%%%%%%%%%%%%%%%%%%%%%%%%%%%%%%%%%%%%%%%%%%%%%%%%
\subsection{Common Errors}
\label{subsec:basic-counting-errors}
%%%%%%%%%%%%%%%%%%%%%%%%%%%%%%%%%%%%%%%%%%%%%%%%%%%%%%%%%%%%

\subsubsection{Adding When Choices Occur in Sequence}

If an outcome requires making one choice and then another, the numbers
of possibilities are generally multiplied.

%%%%%%%%%%%%%%%%%%%%%%%%%%%%%%%%%%%%%%%%%%%%%%%%%%%%%%%%%%%%
\subsubsection{Multiplying Mutually Exclusive Cases}

If an outcome belongs to exactly one of several disjoint alternatives,
the counts should be added rather than multiplied.

%%%%%%%%%%%%%%%%%%%%%%%%%%%%%%%%%%%%%%%%%%%%%%%%%%%%%%%%%%%%
\subsubsection{Ignoring Order}

A committee and an ordered ranking are different objects.

For a committee,

\[
\{A,B,C\}
\]

has no ordering.

For a ranking,

\[
(A,B,C)
\]

and

\[
(B,A,C)
\]

are different.

%%%%%%%%%%%%%%%%%%%%%%%%%%%%%%%%%%%%%%%%%%%%%%%%%%%%%%%%%%%%
\subsubsection{Ignoring Repetition}

A password may allow repeated symbols.

A committee usually does not allow the same person to occupy multiple
member positions.

The model must reflect the actual rules.

%%%%%%%%%%%%%%%%%%%%%%%%%%%%%%%%%%%%%%%%%%%%%%%%%%%%%%%%%%%%
\subsubsection{Using Permutations When Combinations Are Needed}

If order does not matter, using

\[
\frac{n!}{(n-r)!}
\]

counts each subset

\[
r!
\]

times.

The correct unordered count is

\[
\binom{n}{r}.
\]

%%%%%%%%%%%%%%%%%%%%%%%%%%%%%%%%%%%%%%%%%%%%%%%%%%%%%%%%%%%%
\subsubsection{Using Stars and Bars with Distinct Objects}

Stars and bars applies naturally when the distributed objects are
identical and the containers are distinct.

It does not directly count distributions of distinct objects.

%%%%%%%%%%%%%%%%%%%%%%%%%%%%%%%%%%%%%%%%%%%%%%%%%%%%%%%%%%%%
\subsubsection{Ignoring Upper Bounds in Stars and Bars}

The equation

\[
x_1+\cdots+x_k=n
\]

with

\[
x_i\geq0
\]

is different from the same equation with restrictions such as

\[
x_i\leq5.
\]

Upper bounds usually require inclusion--exclusion or another method.

%%%%%%%%%%%%%%%%%%%%%%%%%%%%%%%%%%%%%%%%%%%%%%%%%%%%%%%%%%%%
\subsubsection{Subtracting Intersections Incorrectly}

For two overlapping sets,

\[
|A|+|B|
\]

counts the intersection twice.

One copy must be removed:

\[
|A\cup B|
=
|A|+|B|-|A\cap B|.
\]

For more sets, inclusion--exclusion must continue with alternating
signs.

%%%%%%%%%%%%%%%%%%%%%%%%%%%%%%%%%%%%%%%%%%%%%%%%%%%%%%%%%%%%
\subsubsection{Dividing by a Symmetry Factor Without Justification}

Division works only when every desired object has exactly the same
number of representations.

Do not automatically divide by the size of a symmetry group.

%%%%%%%%%%%%%%%%%%%%%%%%%%%%%%%%%%%%%%%%%%%%%%%%%%%%%%%%%%%%
\subsubsection{Confusing At Least with Exactly}

``At least \(k\)'' includes every case

\[
k,k+1,k+2,\ldots.
\]

``Exactly \(k\)'' includes only one case.

%%%%%%%%%%%%%%%%%%%%%%%%%%%%%%%%%%%%%%%%%%%%%%%%%%%%%%%%%%%%
\subsubsection{Forgetting the Complement Method}

Problems involving phrases such as

\[
\text{at least one}
\]

are often much easier to solve by counting

\[
\text{all}
-
\text{none}.
\]

%%%%%%%%%%%%%%%%%%%%%%%%%%%%%%%%%%%%%%%%%%%%%%%%%%%%%%%%%%%%
\subsection{Applications}
\label{subsec:basic-counting-applications}
%%%%%%%%%%%%%%%%%%%%%%%%%%%%%%%%%%%%%%%%%%%%%%%%%%%%%%%%%%%%

\begin{applicationbox}

\textbf{Probability.}
For equally likely outcomes,

\[
\mathbb{P}(A)
=
\frac{|A|}{|\Omega|},
\]

so probability calculations often reduce directly to counting.

\medskip

\textbf{Computer Science.}
Counting determines the size of search spaces, input spaces, state
spaces, and possible algorithm configurations.

\medskip

\textbf{Cryptography.}
The number of possible passwords, keys, or codes determines the size of
a brute-force search space.

\medskip

\textbf{Scheduling.}
Assignments of employees, machines, aircraft, classrooms, and jobs are
combinatorial configurations.

\medskip

\textbf{Statistics.}
Sampling without replacement uses combinations, while ordered sampling
uses permutations.

\medskip

\textbf{Genetics.}
Possible sequences and inheritance configurations are naturally
enumerative.

\medskip

\textbf{Network Science.}
Graphs, paths, matchings, and connectivity patterns generate large
families of discrete possibilities.

\medskip

\textbf{Operations Research.}
Counting feasible solutions helps reveal the computational difficulty
of assignment, routing, packing, and scheduling problems.

\medskip

\textbf{Coding Theory.}
The number of words at a given Hamming distance is determined using
binomial coefficients.

\end{applicationbox}

%%%%%%%%%%%%%%%%%%%%%%%%%%%%%%%%%%%%%%%%%%%%%%%%%%%%%%%%%%%%
\subsection{Exercises}
\label{subsec:basic-counting-exercises}
%%%%%%%%%%%%%%%%%%%%%%%%%%%%%%%%%%%%%%%%%%%%%%%%%%%%%%%%%%%%

\begin{exercise}[1]
A restaurant offers \(5\) appetizers and \(8\) entrees. If a customer
chooses exactly one item, how many choices are possible?
\end{exercise}

\begin{exercise}[1]
A meal consists of one appetizer and one entree. How many meals are
possible using the choices in the previous exercise?
\end{exercise}

\begin{exercise}[1]
How many binary strings of length \(10\) exist?
\end{exercise}

\begin{exercise}[1]
How many strings of length \(6\) can be formed from an alphabet of
\(4\) symbols if repetition is allowed?
\end{exercise}

\begin{exercise}[1]
How many permutations of \(7\) distinct objects exist?
\end{exercise}

\begin{exercise}[1]
Evaluate

\[
8!.
\]
\end{exercise}

\begin{exercise}[1]
Evaluate

\[
\binom{10}{3}.
\]
\end{exercise}

\begin{exercise}[1]
Evaluate

\[
P(10,3).
\]
\end{exercise}

\begin{exercise}[1]
How many ways can \(3\) students be selected from \(12\) students?
\end{exercise}

\begin{exercise}[1]
How many ways can first, second, and third place be awarded among
\(12\) contestants?
\end{exercise}

\begin{exercise}[1]
How many \(5\)-element subsets does a \(9\)-element set have?
\end{exercise}

\begin{exercise}[1]
How many binary strings of length \(9\) contain exactly four \(1\)'s?
\end{exercise}

\begin{exercise}[1]
How many binary strings of length \(9\) contain no \(1\)'s?
\end{exercise}

\begin{exercise}[1]
How many binary strings of length \(9\) contain at least one \(1\)?
\end{exercise}

\begin{exercise}[1]
How many ways can \(8\) distinct books be arranged on a shelf?
\end{exercise}

\begin{exercise}[1]
How many ways can \(7\) people sit around a circular table?
\end{exercise}

\begin{exercise}[1]
State the pigeonhole principle.
\end{exercise}

\begin{exercise}[1]
If \(31\) objects are placed into \(5\) boxes, show that one box
contains at least \(7\) objects.
\end{exercise}

\begin{exercise}[2]
How many \(6\)-character passwords can be created from \(26\) letters
if repetition is allowed?
\end{exercise}

\begin{exercise}[2]
How many \(6\)-character passwords can be created from \(26\) letters
if repetition is not allowed?
\end{exercise}

\begin{exercise}[2]
How many \(5\)-person committees can be formed from \(15\) people?
\end{exercise}

\begin{exercise}[2]
How many of those committees contain one specified person?
\end{exercise}

\begin{exercise}[2]
How many exclude that specified person?
\end{exercise}

\begin{exercise}[2]
Use the previous two exercises to derive Pascal's identity in this
specific case.
\end{exercise}

\begin{exercise}[2]
How many arrangements of the letters in

\[
\text{MISSISSIPPI}
\]

are distinct?
\end{exercise}

\begin{exercise}[2]
How many nonnegative integer solutions satisfy

\[
x_1+x_2+x_3=15?
\]
\end{exercise}

\begin{exercise}[2]
How many positive integer solutions satisfy

\[
x_1+x_2+x_3=15?
\]
\end{exercise}

\begin{exercise}[2]
How many nonnegative integer solutions satisfy

\[
x_1+x_2+x_3+x_4=20?
\]
\end{exercise}

\begin{exercise}[2]
How many solutions satisfy

\[
x_1+x_2+x_3=20
\]

with

\[
x_1\geq2,
\qquad
x_2\geq4,
\qquad
x_3\geq1?
\]
\end{exercise}

\begin{exercise}[2]
How many multisets of size \(8\) can be formed from \(5\) element
types?
\end{exercise}

\begin{exercise}[2]
How many integers from \(1\) through \(200\) are divisible by \(3\) or
\(5\)?
\end{exercise}

\begin{exercise}[2]
How many integers from \(1\) through \(200\) are divisible by neither
\(3\) nor \(5\)?
\end{exercise}

\begin{exercise}[2]
How many functions exist from a \(5\)-element set to a \(4\)-element
set?
\end{exercise}

\begin{exercise}[2]
How many injective functions exist from a \(4\)-element set to a
\(7\)-element set?
\end{exercise}

\begin{exercise}[2]
How many bijections exist between two \(8\)-element sets?
\end{exercise}

\begin{exercise}[2]
How many monotone lattice paths exist from

\[
(0,0)
\]

to

\[
(6,4)?
\]
\end{exercise}

\begin{exercise}[2]
How many of the paths in the previous exercise pass through

\[
(2,2)?
\]
\end{exercise}

\begin{exercise}[2]
Prove combinatorially that

\[
\binom{n}{r}
=
\binom{n}{n-r}.
\]
\end{exercise}

\begin{exercise}[2]
Prove Pascal's identity combinatorially.
\end{exercise}

\begin{exercise}[2]
Use the binomial theorem to prove

\[
\sum_{k=0}^{n}\binom{n}{k}=2^n.
\]
\end{exercise}

\begin{exercise}[3]
Give a bijective proof that the number of subsets of an \(n\)-element
set is

\[
2^n.
\]
\end{exercise}

\begin{exercise}[3]
Prove Vandermonde's identity combinatorially:

\[
\binom{m+n}{r}
=
\sum_k
\binom{m}{k}
\binom{n}{r-k}.
\]
\end{exercise}

\begin{exercise}[3]
Use double counting to prove

\[
\sum_{k=0}^{n}
k\binom{n}{k}
=
n2^{n-1}.
\]
\end{exercise}

\begin{exercise}[3]
Give a combinatorial proof of

\[
k\binom{n}{k}
=
n\binom{n-1}{k-1}.
\]
\end{exercise}

\begin{exercise}[3]
Prove the handshaking lemma by double counting.
\end{exercise}

\begin{exercise}[3]
Show that every finite graph has an even number of odd-degree vertices.
\end{exercise}

\begin{exercise}[3]
Use inclusion--exclusion to determine how many integers from \(1\)
through \(1000\) are divisible by \(2\), \(3\), or \(5\).
\end{exercise}

\begin{exercise}[3]
Derive the formula

\[
D_n
=
n!
\sum_{k=0}^{n}
\frac{(-1)^k}{k!}
\]

for derangements.
\end{exercise}

\begin{exercise}[3]
Compute

\[
D_5.
\]
\end{exercise}

\begin{exercise}[3]
Use inclusion--exclusion to count surjections from a \(5\)-element set
onto a \(3\)-element set.
\end{exercise}

\begin{exercise}[3]
Show that the number of surjections from an \(n\)-element set onto a
\(k\)-element set equals

\[
k!
\left\{
\begin{matrix}
n\\
k
\end{matrix}
\right\}.
\]
\end{exercise}

\begin{exercise}[3]
Determine the number of nonnegative integer solutions of

\[
x_1+x_2+x_3=12
\]

subject to

\[
x_i\leq5.
\]

Use inclusion--exclusion.
\end{exercise}

\begin{exercise}[3]
How many \(10\)-character binary strings contain exactly four \(1\)'s
and no two \(1\)'s are adjacent?
\end{exercise}

\begin{exercise}[3]
How many ways can \(n\) identical objects be distributed among \(k\)
distinct boxes if exactly \(r\) boxes must be nonempty?
\end{exercise}

\begin{exercise}[3]
Find the number of ways to choose \(8\) objects from \(5\) types if
repetition is allowed and every type must appear at least once.
\end{exercise}

\begin{exercise}[3]
Derive the multinomial coefficient

\[
\frac{n!}{n_1!\cdots n_k!}
\]

using repeated applications of binomial coefficients.
\end{exercise}

\begin{exercise}[4]
Investigate the relationship between binomial coefficients and layers
of the Boolean cube.
\end{exercise}

\begin{exercise}[4]
Study the combinatorial proof of the multinomial theorem.
\end{exercise}

\begin{exercise}[4]
Investigate the number of binary strings of length \(n\) containing no
two consecutive \(1\)'s. Derive a recurrence for the count.
\end{exercise}

\begin{exercise}[4]
Study the number of lattice paths that stay below or on the diagonal

\[
y=x.
\]

Identify the resulting Catalan-number structure.
\end{exercise}

\begin{exercise}[4]
Investigate Stirling numbers of the second kind and derive their
recurrence relation.
\end{exercise}

\begin{exercise}[4]
Investigate Stirling numbers of the first kind and explain their
relationship with permutation cycles.
\end{exercise}

\begin{exercise}[4]
Study Bell numbers and explain how they count set partitions.
\end{exercise}

\begin{exercise}[4]
Investigate the exponential formula connecting set partitions and
exponential generating functions.
\end{exercise}

\begin{exercise}[4]
Study the principle of Möbius inversion on finite posets as a
generalization of inclusion--exclusion.
\end{exercise}

\begin{exercise}[4]
Investigate Burnside's lemma and explain why ordinary division by the
size of a symmetry group does not always correctly count unlabeled
objects.
\end{exercise}

%%%%%%%%%%%%%%%%%%%%%%%%%%%%%%%%%%%%%%%%%%%%%%%%%%%%%%%%%%%%
\subsection{Conceptual Summary}
\label{subsec:basic-counting-summary}
%%%%%%%%%%%%%%%%%%%%%%%%%%%%%%%%%%%%%%%%%%%%%%%%%%%%%%%%%%%%

The addition principle counts disjoint alternatives:

\[
\boxed{
|A_1\cup\cdots\cup A_k|
=
|A_1|+\cdots+|A_k|.
}
\]

The multiplication principle counts sequences of choices:

\[
\boxed{
n_1n_2\cdots n_k.
}
\]

Counting by complement uses

\[
\boxed{
|A|
=
|U|-|A^c|.
}
\]

The pigeonhole principle states that sufficiently many objects placed
into fewer boxes force repetition.

Permutations of \(n\) distinct objects are counted by

\[
\boxed{
n!.
}
\]

Ordered selections without repetition are counted by

\[
\boxed{
\frac{n!}{(n-r)!}.
}
\]

Unordered selections without repetition are counted by

\[
\boxed{
\binom{n}{r}
=
\frac{n!}{r!(n-r)!}.
}
\]

Selections with repetition and order are counted by

\[
\boxed{
n^r.
}
\]

Unordered selections with repetition are counted by

\[
\boxed{
\binom{n+r-1}{r}.
}
\]

Stars and bars gives

\[
\boxed{
x_1+\cdots+x_k=n,
\qquad
x_i\geq0
}
\]

the number of solutions

\[
\boxed{
\binom{n+k-1}{k-1}.
}
\]

Inclusion--exclusion corrects overlapping counts.

For two sets,

\[
\boxed{
|A\cup B|
=
|A|+|B|-|A\cap B|.
}
\]

Derangements satisfy

\[
\boxed{
D_n
=
n!
\sum_{k=0}^{n}
\frac{(-1)^k}{k!}.
}
\]

The number of surjections from an \(m\)-element set onto an
\(n\)-element set is

\[
\boxed{
\sum_{k=0}^{n}
(-1)^k
\binom{n}{k}
(n-k)^m.
}
\]

Double counting proves identities by counting the same finite set in
two different ways.

Bijective proofs establish equalities by showing two families of
objects are structurally equivalent.

These basic principles form the foundation for the more advanced
enumerative methods developed throughout the rest of the chapter.

%%%%%%%%%%%%%%%%%%%%%%%%%%%%%%%%%%%%%%%%%%%%%%%%%%%%%%%%%%%%
\subsection{Section Connections}
\label{subsec:basic-counting-connections}
%%%%%%%%%%%%%%%%%%%%%%%%%%%%%%%%%%%%%%%%%%%%%%%%%%%%%%%%%%%%

\chapterconnection{
Basic Counting builds directly on the finite sets, Cartesian products,
strings, subsets, and discrete configuration spaces introduced in
Section~\ref{sec:discrete-structures}. Recurrence Relations will count
objects by relating structures of size \(n\) to smaller structures.
Generating Functions will encode counting sequences algebraically.
Enumerative Combinatorics will develop systematic techniques involving
partitions, permutations, tableaux, lattice paths, and related
structures. Extremal Combinatorics will replace the question ``how
many?'' with ``how large can a structure be?'' Algebraic Combinatorics
will use symmetry and algebraic structure to count more sophisticated
objects. Probabilistic Combinatorics will combine counting with
probability. Graph Theory, Ramsey Theory, Design Theory, and Matroid
Theory all rely repeatedly on the counting principles established
here.
}

%%%%%%%%%%%%%%%%%%%%%%%%%%%%%%%%%%%%%%%%%%%%%%%%%%%%%%%%%%%%
% SECTION SOURCE TRACKING
%%%%%%%%%%%%%%%%%%%%%%%%%%%%%%%%%%%%%%%%%%%%%%%%%%%%%%%%%%%%

%------------------------------------------------------------
% 6.2 Basic Counting
%------------------------------------------------------------
%
% Source basis:
%   - Standard elementary combinatorics.
%   - Standard discrete mathematics.
%   - Standard finite set counting.
%
% Used for:
%   - addition principle;
%   - multiplication principle;
%   - dependent choices;
%   - decision trees;
%   - complement principle;
%   - division principle;
%   - pigeonhole principle;
%   - generalized pigeonhole principle;
%   - factorials;
%   - permutations;
%   - ordered selections;
%   - falling factorials;
%   - combinations;
%   - binomial coefficients;
%   - Pascal's identity;
%   - Pascal's triangle;
%   - binomial theorem;
%   - binomial coefficient identities;
%   - multinomial coefficients;
%   - repeated-object permutations;
%   - sampling with and without repetition;
%   - stars and bars;
%   - positive and nonnegative integer solutions;
%   - lower and upper bounds;
%   - combinations with repetition;
%   - inclusion--exclusion;
%   - derangements;
%   - counting functions;
%   - injections, surjections, and bijections;
%   - Stirling-number preview;
%   - circular permutations;
%   - constrained subset selection;
%   - counting by cases;
%   - bijective proofs;
%   - double counting;
%   - Vandermonde's identity;
%   - lattice paths;
%   - Boolean-cube layers;
%   - multinomial theorem;
%   - counting by encoding;
%   - overcounting and symmetry cautions;
%   - applications and exercises.
%
% Notes:
%   - Detailed recurrence methods appear in Section 6.3.
%   - Generating functions are developed in Section 6.4.
%   - Stirling, Bell, Catalan, and other classical counting sequences
%     are developed more fully in Enumerative Combinatorics.
%   - Symmetry counting is treated later through algebraic
%     combinatorics and group actions.
%   - Inclusion--exclusion is introduced here in its classical finite
%     form and may later be generalized through Möbius inversion.
%
%%%%%%%%%%%%%%%%%%%%%%%%%%%%%%%%%%%%%%%%%%%%%%%%%%%%%%%%%%%%